% ============================================================
% Section 3: System Architecture
% ============================================================

We begin with a formal characterization of the agent, context window, and
tool set that define the framework, then describe the three-level hierarchy
in detail.

% --------------------------------------------------------
\subsection{Formal Preliminaries}
\label{sec:arch:prelim}
% --------------------------------------------------------

\begin{definition}[Agent]
\label{def:agent}
An \textbf{agent} $A = (M, \tools, \ctx_0, p_{\text{sys}}, k_{\max},
\ell_{\max})$ consists of: a language model $M$ parameterized by weights
$\theta$; a finite tool set $\tools = \{t_1, \ldots, t_m\}$ where each
$t_i : \Sigma^* \to \Sigma^*$ is a deterministic function mapping a JSON
argument string to a string result; an initial context $\ctx_0 \in
\Sigma^*$ (the system prompt); a task-injection prefix $p_{\text{sys}}$; a
maximum tool-call budget $k_{\max} \in \mathbb{N}$; and a maximum answer
token count $\ell_{\max} \in \mathbb{N}$.
\end{definition}

\begin{definition}[Context Window and Pollution Budget]
\label{def:context}
Let $\mathcal{C}_A^{(k)}$ denote the context window of agent $A$ after $k$
tool iterations:
\begin{equation}
  \mathcal{C}_A^{(k)} \;=\; \ctx_0 \;\|\; p_{\text{sys}} \;\|\;
  \bigoplus_{i=1}^{k} \bigl( \text{call}_i \;\|\; \text{obs}_i \bigr),
\end{equation}
where $\|$ denotes string concatenation and $\text{call}_i,\, \text{obs}_i$
are the $i$-th tool call and its observation. The \textbf{pollution budget}
$\rho_A^{(k)} = |\mathcal{C}_A^{(k)}| - |\ctx_0| - |p_{\text{sys}}|$
measures the tokens consumed by tool history after $k$ iterations.
\end{definition}

\noindent Context isolation is achieved by assigning a fresh context
$\mathcal{C}_A^{(0)}$ to every new agent invocation: the pollution budget
is always $\rho_A^{(0)} = 0$ at the start of each worker agent's execution.

% --------------------------------------------------------
\subsection{Three-Level Hierarchical Architecture}
\label{sec:arch:hierarchy}
% --------------------------------------------------------

\sysname implements a three-level hierarchy $(\Hm, \{\Dm_d\}_{d \in
\mathcal{D}}, \{\Wa_i\}_{i \in \mathcal{I}})$ as illustrated in
Figure~\ref{fig:arch}.

\subsubsection{Level 1: High Manager ($\Hm$)}

The singleton High Manager enforces the global four-phase verification
pipeline. It is \emph{not} an LLM-planning agent: its execution graph is
hardcoded as the fixed sequence in Algorithm~\ref{alg:pipeline}, preventing
the class of planning failures (e.g., omitting the verification step) that
afflict unconstrained orchestrators.

\begin{algorithm}[t]
\caption{Level 1 High Manager: Four-Phase Verification Pipeline}
\label{alg:pipeline}
\begin{algorithmic}[1]
\Require Issue $q$, hints $h$, repository $\mathcal{R}$, max retries $k_{\max}$
\Ensure Hierarchical result $r \gets ({\patch}, \textit{verified}, \textit{tokens})$
\State $\text{ctx} \gets \textsc{BuildContext}(q, h, \mathcal{R})$
\State \textbf{// Phase 1: Analysis Domain}
\State $\text{diag} \gets \Dm_{\text{analysis}}.\textsc{Run}(q, \text{ctx})$
\State \textbf{// Phase 2: Testing / Reproduction Domain}
\State $\text{repro} \gets \Dm_{\text{testing}}.\textsc{Run}(q, \text{diag}, \text{ctx})$
\State \textbf{// Phases 3--4: Patch--Verify Loop}
\State $\text{verified} \gets \textbf{false}$; $\;\text{feedback} \gets \varnothing$; $\;k \gets 0$
\While{$k \leq k_{\max}$ \textbf{and not} $\text{verified}$}
  \State $\patch \gets \Dm_{\text{patch}}.\textsc{Run}(q, \text{diag}, \text{repro}, \text{feedback})$
  \State $v \gets \Dm_{\text{testing}}.\textsc{Run}(\text{verify\_goal}, \text{ctx})$
  \If{$\verify(\patch, v) = \texttt{BUG\_FIXED}$}
    \State $\text{verified} \gets \textbf{true}$
  \Else
    \State $\text{feedback} \gets v.\textit{summary}$; $\;k \gets k + 1$
  \EndIf
\EndWhile
\State \Return $(\patch, \text{verified}, \sum_d \text{tokens}(\Dm_d))$
\end{algorithmic}
\end{algorithm}

\subsubsection{Level 2: Domain Managers ($\Dm_d$)}

Each Domain Manager is a full \texttt{SubAgentManager} instance, equipped
with its own LLM-driven planner that dynamically decomposes the domain
goal into subtasks and routes them to Level 3 workers. Formally:

\begin{definition}[Domain Manager]
\label{def:domainmgr}
A \textbf{Domain Manager} $\Dm_d = (\plan_d, \{\Wa_i\}_{i \in d},
s_d, k_d^{\max})$ consists of: a planner LLM $\plan_d$ that produces a
JSON subtask list $\pi = [(\text{id}, \text{task}, \text{agent},
\text{depends\_on})]$; a set of worker agents $\{\Wa_i\}_{i \in d}$ scoped
to domain $d$; an execution strategy $s_d \in \{\text{sequential},
\text{parallel}, \text{adaptive}\}$; and a maximum subtask count $k_d^{\max}$.
\end{definition}

We instantiate three Domain Managers:

\begin{itemize}
  \item $\Dm_{\text{analysis}}$: workers $\{\texttt{issue\_analyzer},
  \texttt{code\_explorer}\}$; strategy $s = \text{sequential}$. Responsible
  for root-cause diagnosis and code localization.

  \item $\Dm_{\text{testing}}$: workers $\{\texttt{reproducer},
  \texttt{test\_generator}, \texttt{test\_runner}\}$; strategy $s =
  \text{sequential}$. Responsible for bug reproduction and patch verification.

  \item $\Dm_{\text{patch}}$: workers $\{\texttt{patch\_writer}\}$;
  strategy $s = \text{sequential}$. Responsible for surgical code
  modification. Contains exactly one worker to prevent conflicting edits.
\end{itemize}

Context passing between Domain Managers flows exclusively through the
\textbf{structured summary interface}: the Level 1 manager receives
$\Dm_d.\textit{summary} \in \Sigma^*$, a concise LLM-synthesized report
(not the raw tool call history), and injects only this summary as context
to the next Domain Manager. This bounded-length summary enforces that
cross-domain information transfer scales as $O(1)$ in the number of tool
calls made within the upstream domain, not $O(k \cdot \ell_{\text{tool}})$.

\subsubsection{Level 3: Worker Agents ($\Wa_i$)}

Each worker agent operates with a minimal, task-specific tool set. The
principle of \textbf{tool minimality}---providing each agent with only the
tools necessary for its specific task---serves two purposes: (i) it reduces
the number of tokens consumed by the tool schema injected into the system
prompt, and (ii) it eliminates failure modes that arise from agents using
tools inappropriately (e.g., \texttt{patch\_writer} using
\texttt{shell\_exec} to run tests rather than applying the patch, exhausting
its iteration budget without modifying any code).

Table~\ref{tab:workers} specifies the tool set, temperature, maximum
iterations, and maximum answer tokens for each worker agent.

\begin{table}[t]
\centering
\caption{Level 3 Worker Agent configurations. $k^{\max}$: max tool
iterations; $\ell^{\max}$: max answer tokens; temp: LLM temperature.}
\label{tab:workers}
\small
\begin{tabular}{llccc}
\toprule
\textbf{Agent} & \textbf{Tools} & $k^{\max}$ & $\ell^{\max}$ & \textbf{Temp} \\
\midrule
\texttt{issue\_analyzer}  & file\_reader, grep, dir\_list & 5  & 2048 & 0.3 \\
\texttt{code\_explorer}   & file\_reader, grep, dir\_list, shell\_exec & 8  & 2048 & 0.2 \\
\texttt{reproducer}       & file\_writer, shell\_exec, view\_file, grep & 5  & 2048 & 0.1 \\
\texttt{test\_generator}  & file\_writer, shell\_exec, view\_file, grep & 5  & 2048 & 0.2 \\
\texttt{patch\_writer}    & str\_replace, view\_file, grep & 6  & 2048 & 0.1 \\
\texttt{test\_runner}     & shell\_exec, file\_reader & 5  & 2048 & 0.1 \\
\bottomrule
\end{tabular}
\end{table}

\paragraph{Mandatory Tool Enforcement.}
The \texttt{reproducer} and \texttt{patch\_writer} agents are configured
with \texttt{mandatory\_tool\_call=True}: the tool loop aborts with an
error if the first model response does not contain a valid tool call.
This prevents the failure mode where the model describes what it \emph{would}
do in text rather than actually calling the tool---a systematic failure
pattern we observed in 23\% of baseline attempts.

\paragraph{Context Window Management.}
Each worker agent's conversation history is managed by a sliding-window
strategy: when the accumulated conversation history exceeds
\texttt{max\_history\_chars} (25,000 chars for \texttt{patch\_writer};
10,000 chars for all others), the middle messages are pruned and replaced
by a one-line summary, always preserving the system prompt and the last
three message pairs. This prevents context exhaustion on large files without
discarding the most recent tool observations.

% --------------------------------------------------------
\subsection{The Skeptical Senior Programmer Verification Loop}
\label{sec:arch:verification}
% --------------------------------------------------------

A fundamental limitation of single-agent systems is that the model that
generated the patch also evaluates whether the patch is correct---a
well-documented failure mode termed \emph{sycophantic self-evaluation}
\cite{huang2023self,sycophancy2023}. We formalize our response to this
problem.

\begin{definition}[Verification Predicate]
\label{def:verify}
Let $\patch \in \Sigma^*$ be a unified diff patch and let $r$ be the output
of executing the reproduction script \texttt{/tmp/reproduce.py} on the
patched repository. The \textbf{verification predicate} is:
\begin{equation}
  \verify(\patch, r) \;=\;
  \begin{cases}
    \texttt{BUG\_FIXED} & \text{if } \texttt{``BUG FIXED''} \in r \;\text{or}\;
    \texttt{``PASSED''} \subseteq r \\
    \texttt{FAIL}       & \text{otherwise.}
  \end{cases}
\end{equation}
\end{definition}

\noindent The critical implementation invariant is that $r$ is always
obtained by the \emph{independent} \texttt{test\_runner} agent, which
receives only the verification goal and the repository path in its context.
The \texttt{test\_runner} has no access to the \texttt{patch\_writer}'s
conversation history and cannot observe what the patch is alleged to do ---
it only observes the actual execution output.

The Level 1 manager enforces a retry budget of $k_{\max} = 2$ patch
attempts. On verification failure, the \texttt{test\_runner}'s output
summary is injected verbatim into the next patch attempt's context as
``Previous Verification Feedback,'' enabling the \texttt{patch\_writer} to
condition on the failure signal without re-executing or re-reading the code.

% --------------------------------------------------------
\subsection{The \texttt{str\_replace} Surgical Edit Paradigm}
\label{sec:arch:strreplce}
% --------------------------------------------------------

A key source of failure in naive LLM-based code editing is full-file
rewriting: the model must faithfully reproduce thousands of characters of
source code from memory, a task where small models hallucinate missing
lines, drop imports, and corrupt indentation. We implement a
\texttt{StrReplaceTool} that accepts three parameters:

\begin{itemize}
  \item \texttt{path}: file path relative to the repository root
  \item \texttt{old\_str}: the exact character sequence to replace
    (whitespace-sensitive, matched verbatim)
  \item \texttt{new\_str}: the replacement character sequence
\end{itemize}

\noindent The tool verifies that \texttt{old\_str} appears \emph{exactly
once} in the file (returning an error if zero or multiple matches are found),
then writes the substituted file. This single-occurrence constraint is
essential: it prevents the model from accidentally replacing the same
fragment in multiple locations, a silent corruption that naive string
replacement would permit.

The \texttt{patch\_writer} system prompt instructs the model to call
\texttt{str\_replace} \emph{before} any file-reading calls, using file
content pre-loaded into its context by the Level 1 manager at plan-time.
This pre-loading strategy---injecting the relevant file's content into the
planner's context and then passing it as \texttt{context} to the
\texttt{patch\_writer} subtask---eliminates the need for
\texttt{patch\_writer} to call \texttt{view\_file} as its first action,
saving one iteration of the $k_{\max} = 6$ budget.

% --------------------------------------------------------
\subsection{Inference Infrastructure}
\label{sec:arch:infra}
% --------------------------------------------------------

All inference is served via a custom \texttt{FastAPI} application wrapping
\texttt{vLLM}'s \texttt{AsyncLLMEngine}~\cite{kwon2023vllm} with the
following configuration:

\begin{itemize}
  \item \textbf{Model}: \texttt{Qwen2.5-Coder-7B-Instruct-AWQ}~\cite{qwen25coder2024},
    7B parameters, 4-bit AWQ quantization~\cite{lin2024awq}.
  \item \textbf{Hardware}: NVIDIA RTX 5080 (16\,GB GDDR7, Blackwell
    architecture). GPU memory utilization: 50\% (8\,GB), reserving 8\,GB
    for the operating system and harness.
  \item \textbf{Context window}: \texttt{max\_model\_len}~$= 16{,}384$ tokens.
  \item \textbf{Attention backend}: FlashAttention~\cite{dao2022flashattention}
    (FlashInfer disabled due to JIT compilation incompatibility with Blackwell).
  \item \textbf{Quantization kernel}: PyTorch Marlin~\cite{frantar2024marlin}
    for NVFP4 MoE layers (Blackwell-optimized path).
  \item \textbf{CPU offload}: 4\,GB model weight offloading with prefetch
    (offload group size 4, 2 in-group, 1-step prefetch) to accommodate the
    7B model's full-precision buffers.
  \item \textbf{Serving protocol}: OpenAI-compatible \texttt{/v1/chat/completions}
    endpoint with real-time streaming; tool calls are parsed from raw model
    text via a dual-format parser that handles both XML
    (\texttt{<tool\_call>}) and JSON block formats.
  \item \textbf{Auto-restart}: a wrapper shell script monitors the vLLM
    process and relaunches it within 30 seconds on crash, providing
    fault tolerance during the 10+ hour benchmark run.
\end{itemize}

\noindent The baseline and hierarchical conditions share the same vLLM
server instance, API base URL, API key, and model name, ensuring that any
observed performance difference is attributable to orchestration architecture
rather than inference infrastructure.
